\scanpage{14}
\mspage{\unsure{84|82}}
dans G. Si C est son centralisateur, alors C est représentable et est
lisse sur S par \inserted{Exp $\overline{\underline{\mathrm{XI}}}$} 5.3., donc la fonction
$s\rightsquigarrow\rho_n(s)=\dim C_s$ est bien
localement constante.

\deleted{Corollaire}
La démonstration de 1.7. est achevée. Nous référerons aux
condition envisagées dans 1.7. b) en disant que dans ce cas, G est
de \uline{rang réductif localement constant}.\deleted{sur S} Notons:

\uline{Corollaire} 1.9. Soit G comme dans 1.7., et soit $s\in S$ tel que
\inserted{$\rho_u(s)=0$ i.e.}
$\rho_r(s)=\rho_{\unsure{m|n}}(s)$, (i.e. les sous-groupes de Cartan de $G_k$ sont des tores, où k
est une clôture algébrique de $k(s)$). Alors il existe un voisinage \deleted{ouvert}
ouvert U de S sur lequel $\rho_r$ et $\rho_n$ sont constants, en particulier pour
tout $t\in U$, le rang unipotent $\rho_u(t)$ de $G_t$ est nul.

En effet, \deleted{\unsure{l'hypothèse s'écrit aussi} $\rho_u(s)=0$, et comme} cela
résulte immédiatement de 1.7. a) et de l'inégalité $\rho_r(t)\le\rho_{\unsure{m|n}}(t)$ pour
tout $t\in S$.

Notons aussi que nous avons prouvé, en même temps que b), le

\uline{Corollaire} 1.10. Soit G comme dans 1.7., et supposons G de rang
réductif localement constant. Considérons le foncteur
\[
\unsure{\mathcal T|\tau}:(\mathrm{Sch})^o\longrightarrow(\mathrm{Ens}),
\]
tel que pour tout S' sur S, on ait
\[
\unsure{\mathcal T|\tau}(S')=\text{ensemble des tores maximaux de }G_{S'}.
\]
Alors $\unsure{\mathcal T|\tau}$ est représentable par un préschéma lisse, séparé et de
type fini sur S.

Il reste à vérifier que $\unsure{\mathcal T|\tau}$ est \deleted{\unsure{séparé}} de type fini sur S\unsure{!|?} La
question étant locale pour la topologie fidèlement plate quasi-compacte,
on peut supposer que G admet un tore maximal T. Mais alors $\unsure{\mathcal T|\tau}$
\deleted{\illegible{1}} est
\deleted{\illegible{1}}\inserted{isomorphe à}
$(G/\uline{\operatorname{Norm}}_G(T)$ (comparer \inserted{Exp $\overline{\underline{\mathrm{XI}}}$} 5.3. bis et 5.5.).
\deleted{\unsure{donc}} Le morphisme
$G\longrightarrow\unsure{\mathcal T|\tau}$ défini par $g\rightsquigarrow\operatorname{int}(g)(T)$ étant surjectif,
et G quasi-compact
sur S, il en est de même de $\unsure{\mathcal T|\tau}$, ce qui achève la démonstration.
\margin{13}
